\section{Microarchitectural Isolation}
\label{sec:hardwaremechanisms}

Microarchitectural isolation is a key primitive enabling Citadel's security.
It can be achieved by partitioning state shared concurrently (spatial sharing) and flushing or sanitizing microarchitectural state shared across time on context switches (temporal sharing).
We use and extend the hardware isolation mechanisms of \cite{bourgeat2019mi6} that we quickly overview here.

\textbf{Spatially Shared Resources}
\label{sec:mi6detailsspacial}
These are the resources shared across cores.
This is essentially the entire shared memory hierarchy starting with the shared Last Level Cache (LLC), including the LLC TLB,  shared Miss Status Handling Registers (MSHRs).
These resources are partitioned such that no partition is ever shared between two security domains.

\textbf{Temporally Shared Resources}
These resources correspond to those private to one core and private partitions of spatially shared resources.
This includes many microarchitectural states used by the out-of-order core such as the execution pipelines, the ROB or the L1 caches.
To achieve isolation, Citadel flushes the L1 caches, the L1 TLB, the L1 replacement policy state, the partitioned MSHRs, the LLC partitions and the LLC arbiters.
To block side channels, Citadel makes sure to sanitize these shared microarchitectural states before context switching between different security domains.
For core-private resources that means the flushing is performed at each context switch, but for LLC partitions we only need to flush when the associated memory region is passed from one security domain to another.

\textbf{Other security considerations}
The SM has access to the entire physical memory as it runs in Machine (M) mode.
That means it is especially vulnerable to speculative attacks and could be exploited to leak any piece of data present on the machine (not only its own secrets but also OS or enclave secrets).
As a result, like in MI6, Citadel's speculation is deactivated for any program running in M mode, preventing the SM from being vulnerable to Spectre-like attacks.